\chapter{Conclusiones y Vías Futuras}
\label{chap:conclusiones}

% \section{Conclusiones}

% \subsection{Logros Principales}

% Este Trabajo de Fin de Grado ha alcanzado exitosamente los objetivos propuestos, demostrando la viabilidad de una plataforma educativa para trading cuantitativo que integra herramientas de producción (Java 21 + Spring Boot) con innovación en ML (Python + multimodelo).

% \subsubsection{Arquitectura Híbrida Robusta}

% Se ha implementado una arquitectura poliglota que separa responsabilidades por idioma y modelo de ejecución:

% \begin{itemize}
%     \item \textbf{Java 21:} Orquestación, persistencia, transaccionalidad, concurrencia con Virtual Threads.
%     \item \textbf{Python:} Cálculo de indicadores, ML, backtesting, ejecución en tiempo real.
% \end{itemize}

% Los beneficios demostrados incluyen:
% \begin{itemize}
%     \item \textbf{Tipo seguro:} Java previene errores en tiempo de compilación que serían costosos en producción.
%     \item \textbf{Agilidad:} Python permite iterar rápidamente en estrategias sin recompilar.
%     \item \textbf{IPC versionado:} El contrato explícito garantiza evolución sin romper clientes.
% \end{itemize}

% \subsubsection{Integridad de Datos Financieros}

% La implementación de transaccionalidad ACID en MySQL, locks pesimistas, y un Ledger inmutable garantiza que:

% \begin{itemize}
%     \item Ningún saldo se modifica sin registro en Ledger.
%     \item Las race conditions entre estrategias concurrentes se previenen con \texttt{LockModeType.PESSIMISTIC\_WRITE}.
%     \item La auditoría es posible en cualquier momento revisando \texttt{LedgerEntry}.
% \end{itemize}

% Esto es fundamental en cualquier software financiero, educativo o no, porque evita la desconfianza del usuario en los resultados.

% \subsubsection{Pipeline ML Profesional}

% Se ha construido un pipeline de ML que aborda problemas reales en trading:

% \begin{itemize}
%     \item \textbf{Desbalance de clases:} uso de métricas ponderadas (F1 weighted).
%     \item \textbf{Lookahead bias:} split temporal respetando causalidad.
%     \item \textbf{Problema de generalización:} multimodelo con validación en datos no vistos.
%     \item \textbf{Reproducibilidad:} almacenamiento de metadata y \texttt{random\_state}.
% \end{itemize}

% Los resultados (XGBoost 64.2\% accuracy) demuestran que, aunque no son de ultra-precisión, son superiores a baselines simples y proporcionan valor educativo.

% \subsubsection{Concurrencia Escalable}

% El uso de Virtual Threads (JDK 21) demuestra que es posible:

% \begin{itemize}
%     \item Ejecutar decenas de estrategias simultáneamente sin agotamiento de recursos.
%     \item Mantener latencia baja (< 300ms end-to-end en tiempo real).
%     \item Escalar sin cambios arquitectónicos radicales.
% \end{itemize}

% \subsubsection{ETL Optimizado}

% El \texttt{FetchService} implementa ETL incremental que logra:

% \begin{itemize}
%     \item Reducción de 96.4\% en tiempo de descarga comparando descarga completa vs incremental (una vez).
%     \item Descarga paralela de chunks sin sobrecarga de API.
%     \item Recuperación automática de huecos históricos sin intervención manual.
% \end{itemize}

% \subsection{Contribuciones Técnicas Diferenciales}

% La solución no es solamente una agregación de tecnologías conocidas, sino que aporta soluciones específicas al dominio educativo del trading cuantitativo:

% \subsubsection{Protocolo IPC Versionado}

% La mayoría de sistemas Python-Java usan approaches ad-hoc (JSON sin estructura, pipes sin sincronización). Este proyecto define un envelope explícito con:

% \begin{itemize}
%     \item Versionado (\texttt{protocol\_version}).
%     \item Framing binario con length-prefixing.
%     \item Separación clara de canales (stdout para IPC, stderr para logs).
%     \item Fallback a legacy JSON para compatibilidad.
% \end{itemize}

% \subsubsection{Feature Engineering Dinámico}

% Cada estrategia define sus propios indicadores y features sin código duplicado. El beneficio es:

% \begin{itemize}
%     \item Usuarios pueden experimentar con distintos espacios de features.
%     \item El pipeline de entrenamiento se adapta dinámicamente.
%     \item Se evita la tentación de overfitting a features fijas.
% \end{itemize}

% \subsubsection{Bloqueos Pesimistas en Operaciones Monetarias}

% Este proyecto usa bloqueos SQL a nivel de transacción:

% \begin{itemize}
%     \item Previene sobregiros incluso con decenas de estrategias simultáneas.
%     \item Es una lección importante de ingeniería de software financiero.
% \end{itemize}

% \subsection{Validación Pedagógica}

% Desde la perspectiva educativa, el sistema demuestra:

% \begin{itemize}
%     \item \textbf{Trading educativo es viable:} paper trading con señales reales de Binance.
%     \item \textbf{Backtesting es reproducible:} indicadores y labels controlados determinísticamente.
%     \item \textbf{ML en finanzas tiene limitaciones:} mostramos que 64\% accuracy es modesto.
%     \item \textbf{Arquitectura predice problemas:} la separación de capas y el IPC versionado hicieron fácil agregar nuevos motores.
% \end{itemize}

% \section{Limitaciones del trabajo}

% \subsection{Alcance Técnico}

% \subsubsection{Modelos de ML No Son Predictivos a Nivel Profesional}

% Los resultados de 64.2\% accuracy son superiores a baseline (50\%) pero insuficientes para trading real. Esto es esperado porque:

% \begin{itemize}
%     \item El mercado tiene componentes aleatorios irreducibles.
%     \item Los indicadores técnicos son heurísticas, no relaciones causales.
%     \item Un dataset de 500-1000 velas es pequeño para ML moderno.
% \end{itemize}

% \subsubsection{Latencia de Tiempo Real (257ms) No Es Ultra-Frecuencia}

% La latencia end-to-end es adecuada para trading en timeframes de 1m o superiores, pero no para scalping de segundos.

% \subsubsection{Cobertura de Tests Incompleta (76\%)}

% Los tests automáticos cubren contratos clave pero no todas las ramas. Faltan tests de:

% \begin{itemize}
%     \item Integración end-to-end en CLI.
%     \item Simulación de timeouts y fallo de API.
%     \item Stress testing con cientos de velas simultáneas.
% \end{itemize}

% \subsection{Alcance Funcional}

% \subsubsection{Falta de Soporte para Trading Real (Solo Paper)}

% El sistema simula trading con dinero ficticio. Esto es por diseño (educativo), pero es una limitación clara.

% \subsubsection{No Hay Análisis de Riesgos Avanzado}

% El sistema calcula Max Drawdown y Win Rate, pero faltan métricas de riesgo como:

% \begin{itemize}
%     \item Value at Risk (VaR).
%     \item Sharpe Ratio y Sortino Ratio.
%     \item Calmar Ratio.
% \end{itemize}

% \subsubsection{Indicadores Técnicos Limitados}

% Se implementan solo 5 indicadores (SMA, EMA, RSI, MACD). Faltan Bollinger Bands, Stochastic, ATR, OBV.

% \subsection{Alcance Operacional}

% \subsubsection{No Hay Monitoreo en Producción}

% El sistema no incluye dashboards en tiempo real, alertas, health checks o métricas Prometheus.

% \subsubsection{Persistencia de Modelos ML Manual}

% Los modelos entrenados se guardan como archivos JSON + PKL. Falta versionado con MLflow.

% \subsubsection{Escalabilidad de BD Limitada a Una Instancia}

% MySQL se asume como instancia única sin replicación. Para escala múltiple, requeriría sharding.

% \subsection{Limitaciones Intrínsecas del Dominio}

% \subsubsection{Datos Históricos No Garantizan Rendimiento Futuro}

% El supuesto fundamental del backtesting es que patrones históricos se repiten. Esto es frecuentemente falso:

% \begin{itemize}
%     \item Cambios regulatorios alteran dinámicas.
%     \item Eventos macroeconómicos son impredecibles.
%     \item Overfitting histórico es fácil y peligroso.
% \end{itemize}

% \section{Vías Futuras}

% \subsection{Corto Plazo (1-2 Meses)}

% \subsubsection{Expandir Indicadores Técnicos}

% Implementar Bollinger Bands, Stochastic, ATR, OBV. 

% \textbf{Esfuerzo:} 2-3 días. \textbf{Riesgo:} bajo.

% \subsubsection{Mejorar Cobertura de Tests}

% Target: cobertura $\geq 90\%$ con tests de integración end-to-end y stress testing.

% \textbf{Esfuerzo:} 1 semana. \textbf{Riesgo:} bajo.

% \subsubsection{Agregar Métricas de Riesgo}

% Implementar Sharpe Ratio, Sortino Ratio, Calmar Ratio para evaluar estrategias profesionalmente.

% \textbf{Esfuerzo:} 3-4 días. \textbf{Riesgo:} bajo.

% \subsection{Mediano Plazo (2-4 Meses)}

% \subsubsection{Dashboard en Tiempo Real (Web UI)}

% Construir interfaz web con React + WebSocket para:

% \begin{itemize}
%     \item Listado de estrategias activas.
%     \item Gráficos de equity curve en vivo.
%     \item Tabla de trades en tiempo real.
%     \item Alertas visuales.
% \end{itemize}

% \textbf{Esfuerzo:} 2-3 semanas. \textbf{Riesgo:} medio (nueva stack frontend).

% \subsubsection{Versionado de Modelos ML (MLflow)}

% Integrar MLflow para registro de entrenamientos, comparación automática y descarga de modelos.

% \textbf{Esfuerzo:} 1 semana. \textbf{Riesgo:} bajo.

% \subsubsection{Soporte para Múltiples Assets}

% Expandir a trading multi-asset simultáneo con correlación y portfolio rebalancing.

% \textbf{Esfuerzo:} 2-3 semanas. \textbf{Riesgo:} medio (cambios arquitectónicos).

% \subsection{Largo Plazo (4+ Meses)}

% \subsubsection{Trading Real (Integración de API Privadas)}

% Habilitar órdenes reales en Binance con gestión de secretos y compliance.

% \textbf{Esfuerzo:} 1-2 meses. \textbf{Riesgo:} alto (dinero real).

% \subsubsection{Deep Learning Avanzado}

% Incorporar LSTM, Transformers, Attention mechanisms y Reinforcement Learning.

% \textbf{Esfuerzo:} 1-2 meses. \textbf{Riesgo:} alto (investigación).

% \subsubsection{Escalabilidad Distribuida}

% Migración hacia múltiples servidores Java, MySQL con replicación, Redis distribuido.

% \textbf{Esfuerzo:} 2-3 meses. \textbf{Riesgo:} alto (complejidad operacional).

% \subsubsection{Marketplace de Estrategias}

% Permitir que usuarios publiquen/compartan estrategias con ranking por rendimiento.

% \textbf{Esfuerzo:} 2-3 meses. \textbf{Riesgo:} medio (regulatorio).

% \section{Reflexión Final}

% Este Trabajo de Fin de Grado ha demostrado que es posible construir una plataforma educativa de trading cuantitativo que sea técnicamente rigurosa, escalable, educativa y mantenible. El sistema no es un producto de trading profesional, pero es un excelente caso de estudio de ingeniería de software aplicado a un dominio claramente acotado y relevante.

% Las vías futuras identificadas son todas técnicamente factibles y agregan incrementalmente más valor. La base arquitectónica es sólida y prepara el terreno para evoluciones futuras sin necesidad de reescrituras masivas.
